%% kapa: exam=1 round=24 year=2013 period=1 form=A subject=economics subject_ko=경제원론 seq=16 number=56
%% source: Q-Net 감정평가사 기출문제 (cst00309) artlSeq=5073802 file=extracted/1008/5073802/2013년 제24회 감정평가사 1차시험 기출문제/1교시(A형).hwp
%% conversion: kordoc 4.24.51 -> tools/gen1.py (auto); review=proofread-sonnet; flags=none
\begin{question}{56}
생수 시장을 양분하고 있는 백두산수와 한라산수의 광고여부에 따른 보수행렬은 아래와 같다. 다음 중 옳은 것을 모두 고른 것은? (단, 각 보수쌍에서 왼쪽은 백두산수의 보수이고, 오른쪽은 한라산수의 보수이다.) \par
{\small\centering
\begin{tabular}{|c|c|c|c|}
\hline
\multirow{4}{*}{\makecell{백두산수}} & \multirow{2}{*}{} & \multicolumn{2}{|c|}{\makecell{한라산수}} \\
\cline{3-4}
 &  & \makecell{광고함} & \makecell{광고안함} \\
\cline{2-4}
 & \makecell{광고함} & \makecell{(25, 15)} & \makecell{(30, 0)} \\
\cline{2-4}
 & \makecell{광고안함} & \makecell{(15, 20)} & \makecell{(40, 5)} \\
\hline
\end{tabular}\par}

\begin{framedtext}
ㄱ. 한라산수는 우월전략을 가지고 있다. \par
ㄴ. 백두산수의 우월전략은 광고를 하는 것이다. \par
ㄷ. 내쉬균형은 모두 광고를 하는 것이다.
\end{framedtext}
\begin{choices}
  \item ㄱ
  \item ㄴ
  \item ㄱ, ㄷ
  \item ㄴ, ㄷ
  \item ㄱ, ㄴ, ㄷ
\end{choices}
\end{question}
\answer{3}
